\documentclass[12pt, a4paper]{article}
  \usepackage[utf8]{inputenc}
  \usepackage[T1]{fontenc}
  \usepackage[french]{babel}
  \usepackage{graphicx}

  \title{La Ville de Mexico}
  \author{}
  \date{}

  \begin{document}

   \maketitle

\section*{Mexico City}

\subsection{General Escobedo}

\subsubsection{Lindo}

\textbf{Mexico}, \textit{capitale} du Mexique, est une \textbf{métropole vibrante et cosmopolite} où s'entremêlent l'ancien et le moderne. Elle compte plus de 21 millions d'habitants, ce qui la place parmi les villes les plus peuplées au monde. Sa richesse historique et culturelle est exceptionnelle : \textbf{ruines aztèques}, temples \textit{aztecs}, palais coloniaux, monuments modernes, musées riches de trésors artistiques, etc. La gastronomie mexicaine, riche et variée, offre une expérience culinaire inoubliable aux visiteurs.

\textbf{un riche héritage musical} qui comprend des styles tels que \textit{la mariachi} et \textit{le salsa}. \textbf{"}
De ce fait j'ai juste adoré tout ça admirer la diversité et l’engouement pour \textbf{le patrimoine culturel} qui sont tant à célébrer. J’ai pu ressentir pleinement cette passion lors de visites aux sites historiques, goûter des spécialités locales savoureuses, et interagir avec les habitants bienveillants\textbf{."}

\textit{Le} réseau de transport ferroviaire mexicain, bien que moins développé que dans d’autres pays, offre néanmoins une possibilité de découvrir \textbf{le} pays d’une manière différente.  Des trajets panoramiques à travers des paysages variés, des montagnes aux déserts, permettent d’apprécier la beauté naturelle de \textit{Mexique} au-delà des grandes villes.  Malheureusement, la couverture du réseau n’est pas optimale et certains trajets peuvent être longs et inconfortables.

En définitive, le Mexique offre une expérience culturelle riche et variée, un voyage inoubliable marqué par la découverte d'un patrimoine musical vibrant et d'une beauté naturelle saisissante, accessible en partie grâce à son réseau ferroviaire, malgré ses imperfections.

Ainsi, si le train ne permet pas une exploration exhaustive du territoire mexicain, il demeure un moyen de transport attrayant pour qui souhaite allier découverte culturelle et contemplation des paysages, contribuant à une expérience de voyage authentique et mémorable.

En France, le réseau ferroviaire est bien plus étendu et performant.

En effet, la comparaison entre les deux réseaux ferroviaires met en lumière les différences d'infrastructures et de développement entre les deux pays.

La SNCF, 

Son \textbf{climat} \textit{varié} avec des \textbf{températures} \textbf{modérées} toute l’\textbf{année} et une \textbf{saison} \textbf{fraîche} en \textbf{hiver}. Cela permet d’apprécier la \textbf{nature} tout au long de l’année. De plus, \textit{Mexico} est reconnue comme le "\textbf{Cœur musical du Mexique}", car elle a été un \textbf{centre} pour les \textbf{musiciens} depuis des \textbf{siècles}. On peut découvrir une \textbf{riche tradition musicale} qui s’est \textbf{perpétuée} et se \textbf{développe} encore au fil du \textbf{temps}.


\begin{figure}[h!]
    \centering
    \includegraphics[width=0.8\textwidth]{figures/mexico_city/general_escobedo/default/fig_1.png}
    \caption{Caption here}
    \label{fig:mexico_city_general_escobedo_1}
\end{figure}


\subsection{B}




Les problemes de performance sont 

les problèmes techniques. La situation est complexe et nécessite une approche multidimensionnelle pour améliorer le fonctionnement urbain. En mettant en place des solutions innovantes à la mobilité douce, on pourrait réduire considérablement les embouteillages dans cette ville de métropoles soudées par son vaste réseau de transport public et ses vastes espaces verts communautaires.


La diversidad cultural y dinámicas urbanas de la Ciudad de México se reflejan también en sus numerosos festivales por cumpleaños que mezclan tradiciones ancestrales con formas artísticas modernas. Por ejemplo, el Festival Internacional del Cervantes acoge actuaciones contemporáneas de danza y actos teatral tradicionales para honrar al poeta Miguel de Cervantes. Las celebraciones por aniversarios urbanos ofrecen un festín sensorial que destaca la herencia viviente de México así como sus tendencias artísticas actuales.
Mexico City prides itself on this blend during its festivals d’anniversaire. The city becomes a kaleidoscope of colors, sounds, and tastes as the population revels in these celebrations that honor both history and modernity. From La Feria de Chapultepec to the Guelaguetza festival at Xochimilco'de los Jazmines y Flor del Monte, each event showcases a unique aspect of Mexico’s cultural identity through traditional music, dance, food, and crafts that have been passed down for generations. These festivals are not just about entertainment; they serve as vital threads in the social fabric, fostering community spirit and pride among residents and visitors alike while reinforcing a collective memory of Mexico’s rich past amidst its dynamic present.

Here is the stylized text:

As dusk falls upon these festivities, \textbf{streetlights} flicker on in synchrony with a vibrant ensemble of mariachi bands. The air is filled with the entrancing melodies that resonate from every corner as people begin to dance freely beneath them, their laughter mixing with music and merriment under the twinkling \textbf{sky}. In Mexico City's energetic atmosphere during these anniversaries celebrations, there seems to be a perpetual motion of joyfulness - locals and visitors alike move in unison as they honor traditions that are intrinsic to their identity. The festivals not only serve as platforms for cultural expression but also actively promote inclusivity among the diverse communities inhabiting this metropolis.

Note: I have used the \textbf{extbf} command to boldify the words "streetlights" and "sky".A

  \end{document}






























































